% Pagina di metodologia/risultati, non legata a un giorno specifico del
% diario. Richiede \usepackage{booktabs} e \usepackage{amsmath} nel
% preambolo di main.tex. Includila con \input{...} dove preferisci nel
% documento (corpo della tesi, non necessariamente nella catena dei giorni).

\section{Benchmark FairMind vs LLM: confronto sui cinque effetti causali}

\subsection*{Obiettivo}

Il benchmark confronta due approcci per il calcolo degli effetti di causal fairness sul dataset Adult Income (UCI): \emph{FairMind}, la libreria del progetto che stima una Bayesian Network sullo Standard Fairness Model (SFM, Ple\v{c}ko e Bareinboim, 2024) e calcola gli effetti tramite inferenza causale formale; e un \emph{Large Language Model} locale (Qwen2.5-7B-Instruct, quantizzato Q4\_K\_M, eseguito via \texttt{llama.cpp}), a cui vengono fornite tabelle di probabilit\`a condizionale pre-aggregate e le formule di identificazione, chiedendogli di applicarle.

\subsection*{Metodologia}

Lo SFM adottato \`e $X \to W \to Y$ con confondente $Z$: attributo protetto $X{=}$\texttt{S2\_gender} ($x_0{=}$Female, $x_1{=}$Male), mediatore $W{=}$\texttt{hours-per-week} (discretizzato in 5 fasce), confondente $Z{=}$\texttt{education}, target $Y{=}$\texttt{T\_income} ($y{=}$\,\texttt{>50K}). I cinque effetti calcolati, secondo le formule di identificazione:
\begin{align*}
\text{TV} &= P(y \mid x_1) - P(y \mid x_0) \\
\text{TE} &= \textstyle\sum_z \big[P(y \mid x_1, z) - P(y \mid x_0, z)\big] \, P(z) \\
\text{SE} &= \text{TV} - \text{TE} \\
\text{DE} &= \textstyle\sum_{z,w} \big[P(y \mid x_1, w, z) - P(y \mid x_0, w, z)\big] \, P(w \mid x_0, z) \, P(z) \\
\text{IE} &= \textstyle\sum_{z,w} P(y \mid x_0, w, z) \, \big[P(w \mid x_1, z) - P(w \mid x_0, z)\big] \, P(z)
\end{align*}
FairMind calcola questi valori tramite inferenza esatta (\emph{VariableElimination} e \emph{CausalInference} di \texttt{pgmpy}) sulla Bayesian Network stimata dai dati. Al LLM, invece, vengono fornite le cinque tabelle di probabilit\`a gi\`a aggregate (anzich\'e i dati grezzi, impraticabili data la finestra di contesto limitata) insieme alle formule sopra, chiedendogli di eseguire il calcolo passo-passo.

\subsection*{Risultati}

\begin{center}
\begin{tabular}{lrrrr}
\toprule
Effetto & FairMind & LLM & Errore assoluto & Errore relativo \\
\midrule
TV & 0.194470 & 0.194500 & 0.000030 & 0.02\% \\
TE & 0.183161 & 0.231167 & 0.048005 & 26.21\% \\
SE & -0.007296 & -0.036667 & 0.029371 & 402.56\% \\
DE & 0.137359 & 0.083167 & 0.054192 & 39.45\% \\
IE & -0.045803 & -0.047000 & 0.001197 & 2.61\% \\
\bottomrule
\end{tabular}
\end{center}

In termini di tempo di calcolo, FairMind completa l'intera pipeline in circa 8 millisecondi, contro i \(\sim\)28 secondi del LLM (\(\sim\)3500 volte pi\`u lento), a fronte di un prompt di circa 6870 token in ingresso e 2037 in uscita.

\subsection*{Analisi}

\begin{itemize}
  \item \textbf{TV e IE} (errore $<3\%$): effetti calcolabili con poche operazioni (una differenza diretta per TV; una somma con fattore stabilizzante per IE, dato che il valore vero \`e prossimo allo zero). Il LLM li approssima molto bene.
  \item \textbf{TE e DE} (errore 26--39\%): richiedono somme su molte combinazioni di $z$ e $(z,w)$ --- rispettivamente 16 e 80 termini --- con pi\`u lookup e moltiplicazioni per termine. L'errore cumulativo del LLM cresce con la complessit\`a dell'aggregazione.
  \item \textbf{SE} (errore 402\%): valore vero prossimo a zero (differenza tra due quantit\`a simili, TV e TE); un piccolo errore assoluto produce un errore relativo enorme. Non \`e tanto un limite del LLM quanto una propriet\`a della metrica stessa.
\end{itemize}

\subsection*{Conclusioni}

FairMind resta il riferimento per accuratezza e velocit\`a: il suo approccio basato su un modello causale formale \`e superiore in entrambe le dimensioni. Un LLM pu\`o fungere da ``calcolatore causale'' affidabile per effetti semplici (TV, IE), ma non per effetti che richiedono aggregazioni multi-variabile estese (TE, DE). Questo supporta l'ipotesi di fondo della tesi: per un'analisi di fairness causale rigorosa, i cinque effetti vanno calcolati con un motore di inferenza formale come FairMind, riservando al LLM il ruolo di interprete e redattore dei risultati gi\`a calcolati.
